\chapter{RC Circuits}
\section{Why RC Circuits Matter}
RC circuits appear as coupling networks, bypass networks, timing circuits, audio
filters, and detector components. They are first-order systems, so their time
behavior is governed by one time constant.

\section{Transient Response}
For a capacitor charging through \(R\) from a DC step \(E\),
\[
v_C(t)=E\left(1-\ee^{-t/(RC)}\right).
\]
For discharge from initial voltage \(V_0\),
\[
v_C(t)=V_0\ee^{-t/(RC)}.
\]
The time constant is
\[
\tau=RC.
\]
After one time constant, a charging capacitor has reached about \(63.2\%\) of
the final value. After about five time constants, it is effectively settled for
many practical purposes.

\section{State-Space Interpretation}
\[
\dot{x}=-\frac{1}{RC}x+\frac{1}{RC}u,
\qquad
x=v_C.
\]
The pole is \(-1/(RC)\). The time constant is the inverse magnitude of that
stable first-order eigenvalue.

\section{Frequency Interpretation}
For output across the capacitor,
\[
G_C(\jj\omega)=\frac{1}{1+\jj\omega RC}.
\]
This is a low-pass response. For output across the resistor,
\[
G_R(\jj\omega)=\frac{\jj\omega RC}{1+\jj\omega RC}.
\]
This is a high-pass response.

\section{Energy Interpretation}
\[
E_C=\frac{1}{2}Cv_C^2.
\]
During charging, energy supplied by the source is partly stored in the electric
field and partly dissipated as heat in the resistor. During discharge, stored
electric-field energy is dissipated in the resistor.

\section{Units Check}
\begin{unitsbox}
\[
RC=\si{\second},
\qquad
\omega_c=\frac{1}{RC}=\si{\per\second},
\qquad
f_c=\frac{1}{2\pi RC}=\si{\hertz}.
\]
\end{unitsbox}

\section{Common Mistakes}
\begin{mistakebox}
A capacitor is an open circuit only for DC steady state. During a transient or
at AC, capacitor current is generally nonzero.
\end{mistakebox}

\section{Worked Example}
Let \(R=\SI{4.7}{\kilo\ohm}\) and \(C=\SI{100}{\nano\farad}\).
\[
\tau=RC=(4.7\times10^3)(100\times10^{-9})
=\SI{4.70e-4}{\second}.
\]
The cutoff frequency is
\[
f_c=\frac{1}{2\pi\tau}\approx \SI{339}{\hertz}.
\]
